% Usage: \begin{paperbox}{Title}{Authors}{Venue & Year}{BibTeX Key}
\begin{paperbox}
  {Example Title: A Clean Framework for Organizing Academic Reading Logs}
  {Doe, J., Smith, A., et al.}
  {Journal of Reusable Science 2026}
  {doe2026example}
  This example demonstrates how to use structured headers, custom callout boxes, and consistent prompts to synthesize academic literature effectively.
\end{paperbox}

\loggap
\begin{itemize}
  \item \textbf{The Problem:} Researchers often take unstructured notes that become difficult to navigate and cite months after reading.
  \item \textbf{Why prior work fails:} Single monolithic documents become too long to compile quickly and lack modular reusability across different writing projects.
\end{itemize}

\logmethod
\begin{itemize}
  \item \textbf{The Clever Trick:} Separating the document styling into a dedicated package file while using modular inputs for individual reading notes.
\end{itemize}

\begin{insightbox}[Core Architectural Principle]
  Keep individual reading notes focused on mathematical mechanisms and actionable takeaways rather than copying abstract summaries.
\end{insightbox}

\logresults
\begin{itemize}
  \item \textbf{Efficiency:} Drastically reduces boilerplate typing by automating dividers, citations, and Table of Contents entries.
\end{itemize}

\loglimitations
\begin{itemize}
  \item \textbf{Stated limits:} Requires familiarity with basic directory structures in LaTeX.
  \item \textbf{My critique:} Works best when paired with an external reference manager like Zotero or a clean BibTeX workflow.
\end{itemize}

\logtakeaways
\begin{itemize}
  \item \textbf{Relevance:} Useful as a standard starting template for literature reviews and graduate qualification exams.
  \item \textbf{Action Item:} Fork this repository on GitHub and replace the sample topic folders with active research domains.
\end{itemize}